\documentclass[11pt]{article}
\usepackage[margin=1in]{geometry}
\usepackage{amsmath,amssymb,amsthm}
\usepackage{graphicx}
\usepackage{hyperref}
\emergencystretch=2em

\newtheorem{theorem}{Theorem}
\newtheorem{lemma}{Lemma}
\newtheorem{corollary}{Corollary}
\theoremstyle{definition}
\newtheorem{definition}{Definition}

\newcommand{\N}{\mathbb{N}}
\newcommand{\R}{\mathbb{R}}
\newcommand{\norm}[1]{\left\lVert #1 \right\rVert}
\newcommand{\supp}{\operatorname{supp}}

\title{A Consolidated Partial Answer to the Smooth Non-Complete Norm Question}
\author{Run fa\_banach\_001, agent lane 16}
\date{June 15, 2026}

\begin{document}
\maketitle

\section*{Status}

This packet records the consolidated likely valid partial results for the question in
Azagra--Dobrowolski--Garcia-Bravo \cite{ADG}.  The result is not a full
solution of the question as stated.  It includes the earlier separable
compact-operator argument and the stronger theorem proving the desired
conclusion for every infinite-dimensional Banach space admitting a bounded
linear injection into some \(c_0(\Gamma)\), hence for every infinite-dimensional
Banach space with a Marku\v{s}evi\v{c} basis.

\section*{Source Question}

The source paper asks, on page 13 immediately before Lemma 2.10:

\begin{quote}
It is not known whether every infinite-dimensional Banach space with a
\(C^1\) equivalent norm possesses a \(C^1\) smooth non-complete norm.
\end{quote}

The same passage notes that for \(C^k\), \(k\geq 2\), the corresponding answer
is positive by work of D'Alessandro and Hajek \cite{DH}.

\begin{center}
\includegraphics[width=0.95\linewidth]{figures/open_question_crop.png}
\end{center}

\section*{Definitions}

All Banach spaces are real. A norm \(p\) on a Banach space \(X\) is
\emph{continuous} if it is continuous for the original Banach norm. It is
\emph{non-complete} if the normed space \((X,p)\) is not complete.

A bounded linear operator \(T:X\to Y\) is \emph{bounded below} if there is
\(c>0\) such that \(\norm{Tx}\geq c\norm{x}\) for every \(x\in X\).

\section*{Idea of the Proof}

In the separable case, countably many norm-one functionals separate points.
After rapidly weighting them, they define a compact injective operator into
\(\ell_2\).  Pulling back the Hilbert norm gives a \(C^\infty\)-smooth norm,
and compactness prevents completeness.

For nonseparable spaces with coordinate structure, \(c_0(\Gamma)\) is the
right replacement.  An infinite-dimensional subspace of \(c_0(\Gamma)\) always
has vectors that can be made almost supported on successively new finite
coordinate sets.  A positive diagonal multiplier can damp those finite sets by
factors tending to zero.  This multiplier remains injective on all of
\(c_0(\Gamma)\), but it is not bounded below on the chosen infinite-dimensional
subspace.

After composing the original injection \(J:X\to c_0(\Gamma)\) with this
diagonal multiplier, the pullback of a \(C^\infty\)-smooth equivalent norm on
\(c_0(\Gamma)\) is a \(C^\infty\)-smooth norm on \(X\).  If that pullback norm
were complete, the open mapping theorem would force it to be equivalent to the
original Banach norm, making the composed operator bounded below.  This
contradicts the diagonal construction.

\begin{theorem}[Separable subcase]
\label{thm:separable}
Let \(X\) be an infinite-dimensional separable Banach space.  Then \(X\)
admits a continuous \(C^\infty\)-smooth non-complete norm.
\end{theorem}

\begin{proof}
Choose a sequence \((u_n)\) dense in the unit sphere of \(X\).  For each
\(n\), pick \(x_n^*\in S_{X^*}\) with \(x_n^*(u_n)=1\).  The family
\((x_n^*)\) separates points: if \(x\neq0\), choose \(n\) with
\(\norm{u_n-x/\norm{x}}<1/2\), and then
\[
    x_n^*(x/\norm{x})\geq x_n^*(u_n)-\norm{u_n-x/\norm{x}}>1/2.
\]

Define
\[
    T:X\longrightarrow \ell_2,\qquad
    Tx=\bigl(2^{-n}x_n^*(x)\bigr)_{n=1}^{\infty}.
\]
Then \(T\) is bounded, injective, and compact, since the finite-rank
truncations \(T_Nx=(2^{-n}x_n^*(x))_{n=1}^N\) converge to \(T\) in operator
norm.  Set \(p(x)=\norm{Tx}_{\ell_2}\).  This is a continuous norm on \(X\),
and it is \(C^\infty\)-smooth on \(X\setminus\{0\}\) because \(T\) is
injective and the Hilbert norm is \(C^\infty\)-smooth off the origin.

If \((X,p)\) were complete, then \(T\) would be an isometry from \((X,p)\)
onto a closed subspace of \(\ell_2\).  Thus \(T:X\to T(X)\), viewed as a
bounded bijection from the original Banach space \(X\) onto the Banach space
\(T(X)\), would have bounded inverse by the open mapping theorem.  Hence \(T\)
would be bounded below in the original norm.  This is impossible for a compact
operator on an infinite-dimensional Banach space.  Therefore \(p\) is
non-complete.
\end{proof}

\begin{lemma}[Diagonal damping in \(c_0(\Gamma)\)]
\label{lem:diagonal}
Let \(Y\) be an infinite-dimensional linear subspace of \(c_0(\Gamma)\).  Then
there is a bounded injective diagonal operator
\[
    D:c_0(\Gamma)\longrightarrow c_0(\Gamma), \qquad (Da)(\gamma)=w_\gamma a(\gamma),
\]
with \(0<w_\gamma\leq 1\), such that \(D|_Y\) is not bounded below.
\end{lemma}

\begin{proof}
We inductively construct \(y_n\in Y\) and pairwise disjoint finite sets
\(F_n\subseteq\Gamma\).  Put \(H_0=\emptyset\).  Having chosen
\(H_{n-1}=F_1\cup\cdots\cup F_{n-1}\), the space
\[
    Y\cap \bigcap_{\gamma\in H_{n-1}}\ker \delta_\gamma
\]
is still infinite-dimensional, because only finitely many coordinate
functionals have been imposed.  Choose \(y_n\) in this space with
\(\norm{y_n}_\infty=1\).  Since \(y_n\in c_0(\Gamma)\), choose a finite
set \(E_n\subseteq\Gamma\) such that
\[
    |y_n(\gamma)|\leq 2^{-n}\qquad(\gamma\notin E_n).
\]
Set \(F_n=E_n\setminus H_{n-1}\) and \(H_n=H_{n-1}\cup F_n\).  Then the
sets \(F_n\) are pairwise disjoint, \(y_n\) vanishes on \(H_{n-1}\), and
\[
    |y_n(\gamma)|\leq 2^{-n}\qquad(\gamma\notin F_n).
\]

Define \(w_\gamma=2^{-n}\) for \(\gamma\in F_n\), and \(w_\gamma=1\) for
\(\gamma\notin\bigcup_n F_n\).  The associated diagonal operator \(D\) is
bounded and injective, since all weights are positive.  For every \(n\),
the coordinates in \(F_n\) are damped by \(2^{-n}\), while all coordinates
outside \(F_n\) already have modulus at most \(2^{-n}\).  Hence
\[
    \norm{Dy_n}_\infty\leq 2^{-n}\qquad(n\in\N).
\]
Thus \(D|_Y\) is not bounded below.
\end{proof}

\begin{theorem}
\label{thm:c0}
Let \(X\) be an infinite-dimensional Banach space.  Suppose that there is a
bounded injective linear operator \(J:X\to c_0(\Gamma)\) for some set
\(\Gamma\).  Then \(X\) admits a continuous \(C^\infty\)-smooth non-complete
norm.
\end{theorem}

\begin{proof}
Let \(Y=JX\).  By Lemma \ref{lem:diagonal}, choose a bounded injective
diagonal operator \(D:c_0(\Gamma)\to c_0(\Gamma)\) such that \(D|_Y\) is not
bounded below.  Then \(DJ:X\to c_0(\Gamma)\) is injective.  It is not bounded
below: if \(\norm{DJx}_\infty\geq c\norm{x}\) for all \(x\in X\), then for
\(y=Jx\in Y\) we would have
\[
    \norm{Dy}_\infty\geq c\norm{x}\geq \frac{c}{\norm{J}}\norm{y}_\infty,
\]
contradicting the choice of \(D\).

Let \(\rho\) be a \(C^\infty\)-smooth equivalent norm on \(c_0(\Gamma)\).
Such norms exist for arbitrary \(\Gamma\), by the \(c_0(\Gamma)\) smooth
renorming theorem of Bible--Smith; see the discussion in
Hajek--Talponen \cite{HT}.  Define
\[
    p(x)=\rho(DJx)\qquad(x\in X).
\]
Since \(DJ\) is injective, \(p\) is a norm.  It is continuous for the original
norm because \(D\), \(J\), and \(\rho\) are continuous.  It is
\(C^\infty\)-smooth on \(X\setminus\{0\}\), because \(DJx\neq 0\) whenever
\(x\neq 0\) and \(\rho\) is \(C^\infty\)-smooth off the origin.

It remains to prove that \(p\) is not complete.  Suppose, to the contrary, that
\((X,p)\) is complete.  The identity map
\[
    I:(X,\norm{\cdot})\longrightarrow (X,p)
\]
is a bounded bijection between Banach spaces, so the open mapping theorem gives
a constant \(C>0\) such that \(\norm{x}\leq C p(x)\) for all \(x\in X\).  Since
\(\rho\) is equivalent to the usual \(c_0(\Gamma)\) norm, there is \(A>0\) with
\(\rho(u)\leq A\norm{u}_\infty\) for all \(u\in c_0(\Gamma)\).  Therefore
\[
    \norm{x}\leq CA\norm{DJx}_\infty\qquad(x\in X),
\]
which says that \(DJ\) is bounded below.  This contradicts the construction of
\(D\).  Hence \(p\) is non-complete.
\end{proof}

\begin{corollary}[M-bases, including separable spaces]
Every infinite-dimensional Banach space with a Marku\v{s}evi\v{c} basis admits
a continuous \(C^\infty\)-smooth non-complete norm.
\end{corollary}

\begin{proof}
It is standard that every Banach space with a Marku\v{s}evi\v{c} basis admits
a continuous linear injection into \(c_0(\Gamma)\); see, for example, the
survey discussion in Hajek--Russo \cite{HRsurvey} and the cited reference
\cite[Theorem II.2.4]{DGZ}.  Theorem \ref{thm:c0} applies.
\end{proof}

\section*{Consequences for the Source Problem}

The source question receives a positive answer for every \(C^1\)-smoothly
renormable infinite-dimensional Banach space that also admits a bounded
injection into \(c_0(\Gamma)\), and in particular for every such space with an
M-basis.  This includes all separable spaces by the classical Marku\v{s}evi\v{c}
theorem, and it also includes many standard nonseparable classes with M-bases
or projectional skeletons.  The direct compact-operator proof of
Theorem \ref{thm:separable} is kept in this consolidated packet because it is
the simplest way to see the separable case.

This does not settle the full \(C^1\) problem.  The proof uses an actual
coordinate injection into \(c_0(\Gamma)\).  A \(C^1\)-smooth equivalent norm
implies strong Asplund-type structure, but the literature search did not
produce a theorem giving an M-basis, a \(c_0(\Gamma)\)-injection, or a quotient
smoothness principle sufficient to run this argument in every such space.

\section*{Literature and Novelty Check}

Searches were run for exact phrases and close variants: ``smooth non-complete
norm'', ``C1 smooth non-complete norm'', ``non-complete norm Banach smooth'',
\(c_0(\Gamma)\) subspaces, separable compact injective operators, and
Marku\v{s}evi\v{c} bases.  The search recovered the source paper, dense smooth
subspace results of Dantas--Hajek--Russo \cite{DHR}, c0-smooth renorming
background \cite{HT}, and M-basis surveys \cite{HRsurvey}, but no exact
published solution of the \(C^1\) source question or of the c0-embeddable
partial theorem above.

The earlier live packet
\texttt{1811.07587\_separable\_smooth\_noncomplete\_norm} has been superseded
by this consolidated packet.  Its separable compact-operator proof is now
included as Theorem \ref{thm:separable}.

\section*{Verification Notes}

The argument is purely functional analytic.  The key points for review are:

\begin{enumerate}
\item Lemma \ref{lem:diagonal} only uses finite-codimensional coordinate
avoidance in the infinite-dimensional subspace \(Y\subseteq c_0(\Gamma)\).
\item The diagonal weights are strictly positive, so \(D\) is injective.
\item The pullback norm is smooth because \(DJ\) is injective.
\item Non-completeness follows from the open mapping theorem and the fact that
\(DJ\) was constructed not to be bounded below.
\item The separable theorem is also independently verified by the compact
injective \(\ell_2\)-operator construction.
\end{enumerate}

\section*{Human Review Recommendation}

Review as the single consolidated packet for this target.  The theorem is not
the full answer to the Azagra--Dobrowolski--Garcia-Bravo question, but it
contains the separable compact-operator proof and the stronger
\(c_0(\Gamma)\)-embeddable theorem, hence the M-basis corollary.

\begin{thebibliography}{9}

\bibitem{ADG}
D. Azagra, T. Dobrowolski, and M. Garcia-Bravo,
\emph{Smooth approximations without critical points of continuous mappings
between Banach spaces, and diffeomorphic extractions of sets},
arXiv:1811.07587.

\bibitem{DH}
S. D'Alessandro and P. Hajek,
\emph{Polynomial algebras and smooth functions in Banach spaces},
J. Funct. Anal. \textbf{266} (2014), 1627--1646.

\bibitem{DGZ}
R. Deville, G. Godefroy, and V. Zizler,
\emph{Smoothness and Renormings in Banach Spaces},
Pitman Monographs and Surveys in Pure and Applied Mathematics, \textbf{64},
Longman, Essex, 1993.

\bibitem{DHR}
S. Dantas, P. Hajek, and T. Russo,
\emph{Smooth and polyhedral norms via fundamental biorthogonal systems},
Int. Math. Res. Not. IMRN \textbf{2023} (2023), 13909--13939;
arXiv:2201.03379.

\bibitem{HRsurvey}
P. Hajek and T. Russo,
\emph{Norming Markushevich bases: recent results and open problems},
arXiv:2402.18442.

\bibitem{HT}
P. Hajek and J. Talponen,
\emph{Some remarks on smooth renormings of Banach spaces},
arXiv:1705.01384.

\end{thebibliography}

\end{document}
